\documentclass[11pt,letterpaper]{article}

\usepackage[spanish,es-noshorthands]{babel}
\usepackage{fontspec}
\setmainfont{Source Serif Pro}
\usepackage{microtype}
\usepackage{geometry}
\usepackage{booktabs}
\usepackage{tabularx}
\usepackage{enumitem}
\usepackage{xcolor}
\usepackage{graphicx}
\usepackage{csquotes}
\usepackage{amsmath}
\usepackage{siunitx}
\usepackage{tikz}
\usetikzlibrary{arrows.meta,positioning,shapes.geometric,fit,backgrounds,patterns}
\usepackage{xurl}
\usepackage[hidelinks]{hyperref}
\usepackage[backend=biber,style=apa,sorting=nyt]{biblatex}
\DeclareLanguageMapping{spanish}{spanish-apa}

\geometry{margin=2.3cm}
% Los rótulos de las figuras usan nodos de ancho fijo con texto alineado a la
% derecha, y la bibliografía compone en modo sloppy: ambos generan avisos de
% línea holgada que son cosméticos y no indican un problema de composición.
% Se silencian para que `make check` siga siendo sensible a los desbordes reales.
\hbadness=10000
\setlength{\emergencystretch}{1em}
\setlength{\parindent}{0pt}
\setlength{\parskip}{0.55em}
\setlist{nosep}
\addbibresource{referencias.bib}
\AtBeginBibliography{\sloppy}

\sisetup{
  output-decimal-marker = {,},
  group-separator = {.},
  group-minimum-digits = 4,
  detect-all
}

% Colores: titular, invasor, remiendo del titular y segunda medición
\definecolor{cIPv4}{RGB}{31,78,121}
\definecolor{cIPv6}{RGB}{116,66,150}
\definecolor{cNAT}{RGB}{191,110,30}
\definecolor{cApnic}{RGB}{40,124,70}

% ---------------------------------------------------------------------------
% Barras de razón sobre eje logarítmico. La paridad (razón igual a 1) está en
% x = 4,909 y la escala es de 4,0909 cm por década, de 10^-1,2 a 10^1.
% ---------------------------------------------------------------------------
\newcommand{\razonbar}[5]{%
  \node[left, text width=4.7cm, align=right, font=\scriptsize] at (-0.22,#1) {#4};%
  \fill[#3] (4.909,#1-0.17) rectangle (#2,#1+0.17);%
  \node[right, font=\scriptsize] at (9.14,#1) {#5};%
}
\newcommand{\notarazon}{Una razón mayor que uno indica que IPv6 es superior en
  esa figura de mérito; menor que uno, que es inferior.}

\title{\vspace{-1.2cm}\textbf{¿Siempre gana la mejor tecnología? El caso de
IPv6}\\
\large Treinta años de superioridad técnica sin sustitución, 1981 a 2026}
\author{Carlos Luis Rojas Aragonés\\
\small Gestión del Desarrollo Tecnológico}
\date{\small 15 de agosto de 2026}

\begin{document}

\maketitle
\vspace{-0.6em}

\section{Introducción}

El avance que analizo es \textbf{IPv6}, la versión 6 del Protocolo de Internet.
Se especificó en diciembre de 1995 \parencite{rfc1883}, se reeditó en diciembre
de 1998 \parencite{rfc2460} y alcanzó la categoría de norma de Internet, STD 86,
en julio de 2017 \parencite{rfc8200}. Su predecesor, IPv4, se publicó en
septiembre de 1981 \parencite{rfc791} y sigue siendo, cuarenta y cinco años
después, el protocolo que todo usuario de Internet necesita.

El 12 de agosto de 2026, la medición de Google situaba el uso de IPv6 en el
\textbf{47,0 por ciento} de sus usuarios \parencite{googlev6} y la de APNIC Labs
en el \textbf{43,2 por ciento} \parencite{apnicv6}. Treinta años después de su
publicación, el sucesor no ha sustituido a nada: convive.

Este caso es exactamente el inverso del que suele estudiarse en gestión
tecnológica. No se trata de un titular excelente que pierde ante un invasor
mediocre pero barato. Se trata de un invasor que es superior en casi todas las
figuras de mérito medibles, que fue diseñado por el mismo organismo que diseñó
al titular, que no tiene competidor alguno y que aun así no consigue desplazarlo.
La pregunta del enunciado, por tanto, admite aquí una respuesta más incómoda que
un simple \enquote{no}: \textbf{la mejor tecnología ni siquiera pierde, se queda
esperando}.

Hay un dato que ordena todo el análisis y conviene adelantar: \textbf{el remiendo
que salvó a IPv4 se publicó antes que IPv6}. La traducción de direcciones de red,
NAT, apareció en mayo de 1994 \parencite{rfc1631} y el direccionamiento sin
clases, CIDR, en septiembre de 1993 \parencite{rfc1519}, mientras que IPv6 no se
especificó hasta diciembre de 1995. Cuando el sucesor llegó, el problema que
justificaba su existencia ya tenía una solución en producción.

\section{Marco de análisis}

Una \textbf{figura de mérito} (FOM) es una medida cuantitativa del desempeño de
una tecnología en la función que presta, independiente de la implementación
concreta que la realiza \parencite{deweck2022}. Cada tecnología recorre en su FOM
una \textbf{curva en S}: arranque lento, mejora rápida y saturación al acercarse
a un límite físico o económico \parencite{foster1986,utterback1994}.

El caso de IPv6 obliga a añadir una segunda familia de conceptos, porque un
protocolo de red no es un producto de consumo individual. Su valor para cada
adoptante depende de cuántos otros lo hayan adoptado ya, que es la definición de
\textbf{externalidad de red} \parencite{katz1985}. Tres consecuencias, todas
documentadas en la literatura y todas visibles en los datos, estructuran el
análisis.

\begin{enumerate}
  \item \textbf{En un mercado con externalidades de red, el FOM que decide no es
    el desempeño del artefacto, sino el número de interlocutores alcanzables.}
    IPv6 nació con ese valor en cero y IPv4 lo tenía en el cien por cien. Ninguna
    ventaja en las demás figuras de mérito compensa esa diferencia inicial
    \parencite{arthur1989,david1985}.
  \item \textbf{Sin compatibilidad hacia atrás no hay adopción incremental.}
    Cuando el adoptante paga el coste completo y no obtiene beneficio hasta que
    el otro extremo también migra, la decisión racional individual es esperar,
    aunque la migración colectiva sea óptima. \textcite{farrell1985} llaman a
    esto \emph{inercia excesiva}, y es el resultado de equilibrio, no una
    patología de gestión.
  \item \textbf{El titular puede mover su propia curva.} La respuesta del
    incumbente no tiene por qué ser adoptar al invasor: puede ser mejorar
    precisamente en el FOM que justificaba al invasor. Es el llamado
    \emph{efecto vela}, un fenómeno discutido en sus casos clásicos
    \parencite{howells2002,mendonca2013}, pero que aquí aparece fechado,
    documentado y con número de norma.
\end{enumerate}

\section{Panorama: dos curvas en S que nunca se cruzan}

La Figura~\ref{fig:curvas} representa el proceso completo. Como FOM del eje
vertical uso la \textbf{fracción de usuarios que puede utilizar el protocolo de
forma nativa}, porque es el único atributo comparable entre las dos transiciones
y porque es, precisamente, el atributo que gobierna las externalidades de red.

\begin{figure}[ht]
  \centering
  \begin{tikzpicture}
    % --- rejilla y eje vertical ---
    \foreach \yv/\lab in {0/{0}, 1.25/{25}, 2.5/{50}, 3.75/{75}, 5.0/{100}} {
      \draw[gray!20] (-0.1,\yv) -- (13.60,\yv);
      \node[left, font=\scriptsize, gray!55!black] at (-0.15,\yv) {\lab\,\%};
    }
    \draw[-{Stealth[length=2mm]}, gray!60!black] (-0.1,-0.05) -- (-0.1,5.55);
    \node[rotate=90, font=\scriptsize, gray!55!black, align=center]
      at (-1.15,2.5) {Usuarios que pueden usar\\el protocolo de forma nativa};
    % --- eje de tiempo con corte ---
    \draw[gray!60!black] (-0.1,0) -- (3.20,0);
    \draw[-{Stealth[length=2mm]}, gray!60!black] (3.60,0) -- (13.60,0);
    \draw[gray!60!black] (3.28,-0.12) -- (3.40,0.12);
    \draw[gray!60!black] (3.36,-0.12) -- (3.48,0.12);
    \foreach \xv/\lab in {0.237/1983, 1.067/1990, 1.659/1995, 2.252/2000,
                          2.844/2005} {
      \draw[gray!60!black] (\xv,0) -- (\xv,-0.13);
      \node[below, font=\tiny, gray!55!black] at (\xv,-0.13) {\lab};
    }
    \foreach \xv/\lab in {4.213/2010, 5.744/2015, 7.275/2020, 8.806/2025,
                          10.338/2030, 11.869/2035, 13.400/2040} {
      \draw[gray!60!black] (\xv,0) -- (\xv,-0.13);
      \node[below, font=\tiny, gray!55!black] at (\xv,-0.13) {\lab};
    }
    \node[font=\tiny, gray!55!black, align=center] at (1.75,-0.62)
      {escala comprimida};
    \node[font=\tiny, gray!55!black, align=center] at (11.60,-0.62)
      {proyección};
    % --- IPv4 ---
    \draw[very thick, cIPv4] (0,0) -- (0.178,0.50) -- (0.296,5.0) -- (3.20,5.0);
    \draw[very thick, cIPv4] (3.60,5.0) -- (13.60,5.0);
    \draw[cIPv4!70, thin] (0.296,4.55) -- (0.80,3.55);
    \node[font=\tiny, cIPv4, align=left, right] at (0.72,3.30)
      {1 de enero de 1983: unos 400 servidores\\[-0.25em]
       cambian de protocolo en un solo día};
    \node[font=\scriptsize, cIPv4, below] at (11.35,4.96)
      {IPv4: 100\,\% desde 1983, sin descenso};
    % --- publicación de IPv6 ---
    \draw[cIPv6!70, densely dotted, thick] (1.659,0) -- (1.659,1.45);
    \node[font=\tiny, cIPv6, align=center, above] at (1.659,1.45)
      {RFC 1883\\[-0.3em] dic. 1995};
    % --- IPv6, medición de Google ---
    \draw[very thick, cIPv6] (3.6000,0.0082) -- (3.9062,0.0110)
      -- (4.2125,0.0122) -- (4.5187,0.0157) -- (4.8250,0.0342)
      -- (5.1312,0.0778) -- (5.4375,0.1798) -- (5.7438,0.3348)
      -- (6.0500,0.5492) -- (6.3563,0.8406) -- (6.6625,1.0597)
      -- (6.9688,1.2744) -- (7.2750,1.4971) -- (7.5813,1.6668)
      -- (7.8875,1.8807) -- (8.1937,2.0509) -- (8.5000,2.1714)
      -- (8.8063,2.2920) -- (9.1125,2.3561);
    \foreach \x/\y in {5.7438/0.3348, 6.3563/0.8406, 6.9688/1.2744,
                       7.5813/1.6668, 8.1937/2.0509, 8.8063/2.2920,
                       9.1125/2.3561} {
      \fill[cIPv6] (\x,\y) circle (1.7pt);
    }
    % --- IPv6, medición de APNIC ---
    \draw[thick, cApnic, dashed] (5.1312,0.0699) -- (5.4375,0.1108)
      -- (5.7438,0.2055) -- (6.0500,0.3239) -- (6.3563,0.6741)
      -- (6.6625,0.8723) -- (6.9688,1.0847) -- (7.2750,1.2856)
      -- (7.5813,1.3995) -- (7.8875,1.5531) -- (8.1937,1.7432)
      -- (8.5000,1.9010) -- (8.8063,2.0385) -- (9.1125,2.1487);
    \fill[cApnic] (9.1125,2.1487) circle (1.7pt);
    \node[font=\tiny, cApnic, below right=-0.4mm and 0.3mm] at (9.1125,2.1487)
      {43,0\,\%};
    \node[font=\tiny, cIPv6, above left=-0.5mm and -0.5mm] at (9.1125,2.3561)
      {47,1\,\%};
    % --- proyecciones ---
    \draw[thick, cIPv6, densely dashed] (9.1125,2.3561) -- (13.40,4.3379);
    \draw[-{Stealth[length=2mm]}, thick, cIPv6, densely dashed]
      (13.40,4.3379) -- (13.60,4.43);
    \node[font=\tiny, cIPv6, rotate=27, above] at (11.60,3.30)
      {extrapolación lineal: 100\,\% en 2045};
    \draw[cIPv6!55, dotted, thick] (9.1125,2.3775) -- (13.55,2.3775);
    \node[font=\tiny, cIPv6!75!black, below right] at (10.60,2.36)
      {asíntota del ajuste logístico: 47,6\,\%};
    % --- umbral del 50% ---
    \draw[red!65!black, thick, densely dotted] (9.05,2.5) -- (13.55,2.5);
    \draw[red!65!black, thick] (9.186,2.5) circle (3.0pt);
    \draw[red!65!black, thin] (9.10,2.66) -- (8.55,3.30);
    \node[font=\scriptsize, red!65!black, align=center, above left=-0.5mm and -1mm]
      at (8.55,3.30) {28 de marzo de 2026:\\[-0.3em] \tiny primer día por encima
      del 50\,\%};
    % --- leyenda ---
    \begin{scope}[shift={(3.90,5.38)}]
      \draw[very thick, cIPv4] (0,0) -- (0.45,0);
      \node[right, font=\tiny] at (0.5,0) {IPv4};
      \draw[very thick, cIPv6] (1.55,0) -- (2.00,0);
      \node[right, font=\tiny] at (2.05,0) {IPv6, medición de Google};
      \draw[thick, cApnic, dashed] (5.30,0) -- (5.75,0);
      \node[right, font=\tiny] at (5.80,0) {IPv6, medición de APNIC Labs};
    \end{scope}
  \end{tikzpicture}
  \caption{Curvas en S de las dos transiciones de protocolo de Internet sobre un
  mismo FOM. La curva de IPv4 recoge el cambio de NCP a TCP/IP planificado en
  \textcite{rfc801} y ejecutado el 1 de enero de 1983 sobre unos 400 servidores.
  La de IPv6 son las medias anuales de las dos series públicas de medición: la
  de Google, diaria desde el 4 de septiembre de 2008 \parencite{googlev6}, y la
  de APNIC Labs, suavizada a 30 días desde el 7 de octubre de 2013
  \parencite{apnicv6}, ambas descargadas el 15 de agosto de 2026. La
  extrapolación lineal es la regresión de los últimos cinco años de la serie de
  Google, con pendiente de 2,73 puntos porcentuales por año. La asíntota
  logística procede de ajustar $L/(1+e^{-k(t-t_0)})$ a los \num{6552} datos
  diarios por mínimos cuadrados. Ambas extrapolaciones se discuten en la
  sección de precisiones y limitaciones. Elaboración propia.}
  \label{fig:curvas}
\end{figure}

Tres lecturas inmediatas de la figura, que las secciones siguientes desarrollan.

\textbf{Primera: el cruce no ha ocurrido, y la comparación de las dos
transiciones es de otro orden.} El paso de NCP a TCP/IP se planificó en noviembre
de 1981 \parencite{rfc801} y se ejecutó el 1 de enero de 1983, en un día
señalado, sobre unos 400 servidores. La transición a IPv6 lleva desde 1995 y
alcanza al 47 por ciento de unos \num{6100} millones de usuarios
\parencite{itu2025}. La diferencia entre las dos curvas no mide la calidad de los
protocolos: mide el tamaño de la red que hay que coordinar.

\textbf{Segunda: la curva del titular nunca baja.} Esta es la anomalía central
del caso. En una sustitución tecnológica normal, la curva del invasor sube
mientras la del titular cae. Aquí IPv6 sube y IPv4 sigue en el cien por cien,
porque el 47 por ciento que usa IPv6 lo hace en \emph{doble pila}, manteniendo
IPv4 en paralelo, o bien atravesando un traductor hacia IPv4. En ambos casos IPv4
sigue haciendo falta. No hay sustitución, hay adición, y la adición cuesta dinero
en lugar de ahorrarlo.

\textbf{Tercera: la curva de IPv6 ya perdió velocidad.} El incremento anual
máximo se alcanzó entre 2016 y 2017, con 5,83 puntos porcentuales, y en los dos
últimos años ha sido de 2,41 puntos. Un ajuste logístico a la serie completa de
Google sitúa el punto de inflexión en abril de 2019 y la asíntota en el 47,6 por
ciento; el mismo ajuste sobre la serie independiente de APNIC Labs sitúa la
inflexión en octubre de 2019 y la asíntota en el 43,3 por ciento. Es decir, dos
sistemas de medición distintos, con metodologías distintas, coinciden en que la
curvatura ya es negativa.

\section{La tecnología establecida: IPv4 con traducción de direcciones}

\subsection*{01. El desempeño que la hizo dominante}

IPv4 define una dirección de 32 bits, lo que produce un espacio de
\num{4294967296} direcciones \parencite{rfc791}. Con unos \num{6100} millones de
usuarios de Internet en 2026 \parencite{itu2025}, eso equivale a 0,70 direcciones
por usuario antes de contar servidores, routers, sensores ni teléfonos. El
espacio es, en sentido estricto, insuficiente, y lo es desde hace más de una
década: IANA agotó su reserva central el 3 de febrero de 2011 y las cinco
autoridades regionales fueron agotando la suya a partir de esa fecha, empezando
por APNIC en abril de 2011 \parencite{nro2026}.

Lo relevante para la gestión tecnológica no es que se agotara, sino lo poco que
pasó cuando se agotó. La razón es que el FOM que decide no era el número de
direcciones asignables, sino el número de dispositivos que se pueden conectar, y
esos dos números dejaron de ser el mismo en 1994.

\subsection*{02. El titular que mejora en el FOM del invasor}

La secuencia de remiendos de IPv4 está fechada y numerada, y es el argumento más
fuerte del caso:

\begin{itemize}
  \item \textbf{Septiembre de 1993}: CIDR elimina las clases de direcciones y
    permite asignar bloques del tamaño realmente necesario \parencite{rfc1519}.
  \item \textbf{Mayo de 1994}: NAT permite que muchas máquinas compartan una sola
    dirección pública multiplexando puertos \parencite{rfc1631}, formalizado
    después en \textcite{rfc3022}.
  \item \textbf{Febrero de 1996}: RFC 1918 reserva tres bloques privados, de modo
    que el direccionamiento interno deja de consumir espacio público
    \parencite{rfc1918}.
  \item \textbf{Abril de 2012}: RFC 6598 reserva el bloque \texttt{100.64.0.0/10}
    para NAT de operador, la CGNAT, que lleva el mismo truco al nivel del
    proveedor de acceso \parencite{rfc6598}.
\end{itemize}

Las dos primeras son anteriores a IPv6 y la tercera es prácticamente simultánea.
El efecto agregado es que IPv4 dejó de tener un techo duro de direcciones: un
operador que asigna una dirección pública por cada 64 abonados multiplica por 64
su capacidad, y la penalización principal es el número de puertos disponibles por
abonado, que cae de unos \num{64512} a unos \num{1024}
\parencite{rfc7021,ipv4global2025}.

\begin{figure}[ht]
  \centering
  \begin{tikzpicture}
    % ---------------- panel A: cronología comparada ----------------
    \node[font=\scriptsize\itshape, gray!35!black, right] at (-0.1,2.55)
      {Panel A. La carrera entre el sucesor y el remiendo};
    \draw[gray!60!black] (0,0.75) -- (13.30,0.75);
    \foreach \xv/\lab in {0/1990, 2.321/1995, 4.643/2000, 6.964/2005,
                          9.286/2010, 11.607/2015} {
      \draw[gray!60!black] (\xv,0.75) -- (\xv,0.62);
      \node[below, font=\tiny, gray!55!black] at (\xv,0.62) {\lab};
    }
    % --- pista superior: IPv6 ---
    \foreach \xv/\lab in {2.763/{RFC 1883\\IPv6}, 4.156/{RFC 2460},
                          10.415/{World IPv6\\Launch}, 12.771/{RFC 8200\\STD 86}} {
      \draw[cIPv6, thick] (\xv,0.75) -- (\xv,1.15);
      \fill[cIPv6] (\xv,1.15) circle (1.7pt);
      \node[above, font=\tiny, cIPv6, align=center] at (\xv,1.20) {\lab};
    }
    % --- pista inferior: remiendos de IPv4, a dos profundidades ---
    \foreach \xv/\lab in {2.833/{RFC 1918\\privadas},
                          5.108/{RFC 3022}, 10.354/{RFC 6598\\CGNAT}} {
      \draw[cNAT, thick] (\xv,0.75) -- (\xv,0.42);
      \fill[cNAT] (\xv,0.42) circle (1.7pt);
      \node[below, font=\tiny, cNAT, align=center] at (\xv,0.37) {\lab};
    }
    \draw[cNAT, thick] (1.718,0.75) -- (1.718,0.42);
    \fill[cNAT] (1.718,0.42) circle (1.7pt);
    \node[anchor=north east, font=\tiny, cNAT, align=right] at (1.80,0.37)
      {RFC 1519\\CIDR};
    \draw[cNAT, very thick] (2.042,0.75) -- (2.042,-0.34);
    \fill[cNAT] (2.042,-0.34) circle (2.2pt);
    \node[below, font=\tiny, cNAT, align=center] at (2.042,-0.39)
      {RFC 1631\\\textbf{NAT}};
    \draw[gray!55, thick] (9.797,0.75) -- (9.797,2.00);
    \fill[gray!55!black] (9.797,2.00) circle (1.7pt);
    \node[above, font=\tiny, gray!35!black, align=center] at (9.797,2.05)
      {agotamiento de IANA};
    % --- llave del adelanto ---
    \draw[red!45, dotted] (2.042,0.78) -- (2.042,2.05);
    \draw[red!45, dotted] (2.763,1.66) -- (2.763,2.05);
    \draw[{Stealth[length=1.6mm]}-{Stealth[length=1.6mm]}, red!65!black, thick]
      (2.042,2.10) -- (2.763,2.10);
    \node[font=\tiny, red!65!black, align=left, right] at (2.86,2.10)
      {NAT se publicó \textbf{diecinueve meses antes} que IPv6};

    % ---------------- panel B: precio de una dirección IPv4 ----------
    \begin{scope}[shift={(0,-4.55)}]
      \node[font=\scriptsize\itshape, gray!35!black, right] at (-0.1,3.00)
        {Panel B. Precio de mercado de una dirección IPv4, en dólares};
      \foreach \yv/\lab in {0/0, 0.867/20, 1.733/40, 2.600/60} {
        \draw[gray!20] (0,\yv) -- (13.30,\yv);
        \node[left, font=\tiny, gray!55!black] at (-0.05,\yv) {\lab};
      }
      \draw[gray!60!black] (0,0) -- (13.30,0);
      \foreach \xv/\lab in {1.857/2021, 3.714/2022, 5.571/2023, 7.429/2024,
                            9.286/2025, 11.143/2026} {
        \draw[gray!60!black] (\xv,0) -- (\xv,-0.13);
        \node[below, font=\tiny, gray!55!black] at (\xv,-0.13) {\lab};
      }
      \fill[cNAT!14] (1.764,1.170) -- (3.621,2.383) -- (5.571,2.253)
        -- (9.471,1.430) -- (11.050,0.953) -- (11.886,0.848)
        -- (11.886,0) -- (1.764,0) -- cycle;
      \draw[very thick, cNAT] (1.764,1.170) -- (3.621,2.383) -- (5.571,2.253)
        -- (9.471,1.430) -- (11.050,0.953) -- (11.886,0.848);
      \foreach \x/\y in {1.764/1.170, 3.621/2.383, 5.571/2.253, 9.471/1.430,
                         11.050/0.953, 11.886/0.848} {
        \fill[cNAT] (\x,\y) circle (1.7pt);
      }
      \node[font=\tiny, cNAT, above] at (3.621,2.44) {máximo: 55};
      \node[font=\tiny, cNAT, right] at (11.98,0.848) {19,57};
      \node[font=\scriptsize, red!65!black, align=center]
        at (9.55,2.32) {la escasez se relaja,\\[-0.3em] no se agrava};
      \draw[-{Stealth[length=2mm]}, thick, red!65!black]
        (10.35,1.98) -- (11.30,1.05);
    \end{scope}
  \end{tikzpicture}
  \caption{El contrainvasor. Panel A: fechas de publicación de las normas del
  IETF. Las de la pista inferior mejoran a IPv4 en la única figura de mérito que
  justificaba a IPv6, y dos de ellas son anteriores a la especificación del
  sucesor. Panel B: precio medio de una dirección IPv4 en el mercado secundario.
  Los valores de 2020 a 2025 proceden del análisis anual de
  \textcite{huston2026a}, que informa de máximos de 45 a 60 dólares en 2023, de
  26 a 40 a comienzos de 2025 y de una media de unos 22 dólares al cierre de
  2025; los de 2026 proceden de los informes mensuales elaborados sobre los datos
  de transacciones de IPv4.Global \parencite{ipv4center2026,circleid2026}. Es un
  mercado opaco, con precios muy dispersos según el tamaño del bloque, hasta el
  punto de que los bloques de \texttt{/16} o mayores cotizaban por debajo de 10
  dólares en febrero de 2026: los valores deben leerse como orden de magnitud y
  tendencia, no como una serie de precio único. Elaboración propia.}
  \label{fig:contra}
\end{figure}

El Panel B de la Figura~\ref{fig:contra} cierra el argumento económico. Si la
escasez de direcciones IPv4 fuera creciente, el precio subiría y en algún momento
haría racional migrar. Lo que muestran los datos de mercado es lo contrario: el
precio tocó techo alrededor de los 55 dólares por dirección a finales de 2021 y
ha caído desde entonces hasta unos 20 dólares en 2026, un descenso interanual del
39 por ciento en enero de ese año \parencite{ipv4center2026}. El único motor
económico que quedaba para la migración se está apagando.

\section{El invasor: IPv6}

\subsection*{01. En qué es superior, y por cuánto}

IPv6 gana en casi todo lo que se puede medir. El espacio de direcciones pasa de
32 a 128 bits, de \num{4294967296} a $3{,}40\times10^{38}$, un factor de
$2^{96}$, es decir $7{,}92\times10^{28}$. La cabecera pasa de doce campos de
longitud variable a ocho campos de longitud fija de 40 octetos, elimina la suma
de comprobación que IPv4 obliga a recalcular en cada salto y prohíbe que los
routers intermedios fragmenten paquetes \parencite{rfc791,rfc8200}. Añade
autoconfiguración sin servidor \parencite{rfc4862}. Y la agregación funciona:
el 15 de agosto de 2026 la tabla de rutas global de IPv4 tenía \num{1116753}
entradas frente a \num{260375} de IPv6 \parencite{potaroo2026}, es decir IPv6
transporta $7{,}9\times10^{28}$ veces más espacio de direcciones con una cuarta parte de las
entradas de encaminamiento.

\subsection*{02. En qué no es superior, y por qué eso lo decidió todo}

\begin{figure}[ht]
  \centering
  \begin{tikzpicture}
    % ---------------- panel A: espacio de direcciones ----------------
    \node[font=\scriptsize\itshape, gray!35!black, right] at (-4.95,1.35)
      {Panel A. La figura de mérito que no cabe en el mismo eje};
    \foreach \e in {0,5,...,30} {
      \pgfmathsetmacro{\xg}{0.30*\e}
      \draw[gray!25] (\xg,-0.40) -- (\xg,0.90);
      \node[below, font=\tiny, gray!55!black] at (\xg,-0.40) {$10^{\e}$};
    }
    \node[left, text width=4.7cm, align=right, font=\scriptsize] at (-0.22,0.32)
      {Espacio de direcciones\\\tiny (razón entre IPv6 e IPv4)};
    \fill[cIPv6] (0,0.15) rectangle (8.670,0.49);
    \node[right, font=\scriptsize] at (8.77,0.32) {$7{,}92\times10^{28}$};

    % ---------------- panel B: las demás ----------------
    \begin{scope}[shift={(0,-5.35)}]
      \node[font=\scriptsize\itshape, gray!35!black, right] at (-4.95,3.55)
        {Panel B. Las figuras de mérito que sí caben, y que decidieron el
         resultado};
      \foreach \xv/\lab in {0.818/{0,1}, 2.446/{0,25}, 3.677/{0,5},
                            4.909/{1}, 6.140/{2}, 7.771/{5}, 9.000/{10}} {
        \draw[gray!22] (\xv,-0.50) -- (\xv,3.15);
        \node[below, font=\tiny, gray!55!black] at (\xv,-0.50) {\lab};
      }
      \draw[gray!60!black, thick] (4.909,-0.50) -- (4.909,3.15);
      \node[above, font=\tiny, gray!45!black] at (4.909,3.17) {paridad};
      \razonbar{2.85}{7.494}{cIPv6}{Tabla de rutas global\\\tiny
        (\num{1116753} frente a \num{260375} entradas; menor es mejor)}{4,29}
      \razonbar{1.90}{5.629}{cIPv6}{Campos de la cabecera\\\tiny (12 frente a 8;
        menor es mejor)}{1,50}
      \razonbar{0.95}{4.840}{cIPv6!45}{Tiempo de ida y vuelta\\\tiny (medido en
        redes de doble pila nativo)}{$\approx$ 1,0}
      \razonbar{0.00}{3.677}{cNAT}{Sobrecarga por paquete\\\tiny (20 frente a 40
        octetos de cabecera; menor es mejor)}{0,50}
      \node[left, text width=4.7cm, align=right, font=\scriptsize]
        at (-0.22,-1.15) {Interlocutores alcanzables en 1995\\\tiny (el FOM que
        gobierna la adopción)};
      \draw[-{Stealth[length=2mm]}, very thick, red!65!black]
        (4.909,-1.15) -- (0.30,-1.15);
      \node[right, font=\scriptsize, red!65!black] at (5.05,-1.15)
        {$\rightarrow$ 0: fuera de escala por la izquierda};
    \end{scope}
  \end{tikzpicture}
  \caption{Perfil de figuras de mérito de IPv6 expresado como razón frente a
  IPv4, sobre ejes logarítmicos. \notarazon{} Las razones se han orientado de
  modo que un valor mayor que uno signifique siempre ventaja de IPv6. Panel A: la
  ventaja en espacio de direcciones es de 29 órdenes de magnitud y no cabe en el
  mismo eje que las demás. Panel B: el resto de las figuras de mérito se mueve
  entre 0,5 y 4,3, y en dos de ellas IPv6 empata o pierde. La sobrecarga de
  cabecera es peor por diseño, porque la cabecera fija de 40 octetos duplica los
  20 octetos mínimos de IPv4 \parencite{rfc791,rfc8200}. El tiempo de ida y
  vuelta se sitúa en la paridad: \textcite{howard2019} encuentra que las
  operadoras de doble pila nativo miden esencialmente la misma velocidad en los
  dos protocolos. La última fila no es una razón medida sino el valor que tomaba
  el FOM decisivo el día en que se publicó IPv6. Elaboración propia.}
  \label{fig:razon}
\end{figure}

La Figura~\ref{fig:razon} contiene, a mi juicio, la explicación completa del
caso. Leída de arriba a abajo dice tres cosas.

\textbf{La ventaja de IPv6 está concentrada en una sola figura de mérito.} Es
enorme, 29 órdenes de magnitud, pero es una sola, y es además la única que el
usuario final no puede percibir. Nadie nota que su proveedor tiene direcciones de
sobra.

\textbf{En las demás, IPv6 empata o pierde.} La cabecera fija de 40 octetos
duplica los 20 octetos mínimos de IPv4, de modo que en enlaces con paquetes
pequeños IPv6 desperdicia más ancho de banda. En velocidad no hay diferencia
apreciable en redes de doble pila nativo \parencite{howard2019}. Geoff Huston, el
científico jefe de APNIC, lo resume sin rodeos: IPv6 \enquote{no era más rápido,
ni más versátil, ni más seguro que IPv4} \parencite{huston2024}. Su único
beneficio real era mitigar un riesgo futuro, y los riesgos futuros no se
presupuestan.

\textbf{La figura de mérito que decidió el resultado valía cero.} IPv6 no es
compatible hacia atrás con IPv4: una máquina que solo hable IPv6 no puede
comunicarse con una que solo hable IPv4 sin un traductor intermedio. Esto no es
una consecuencia imprevista, sino una decisión de diseño que el propio organismo
reconoció después. En un panel del IETF en San Francisco en marzo de 2009, Leslie
Daigle, entonces directora de tecnología de Internet de la Internet Society,
declaró que \enquote{la falta de compatibilidad real hacia atrás con IPv4 fue el
único fallo crítico} \parencite{daigle2009}.

\subsection*{03. La consecuencia: nadie migra, todos suman}

La estrategia de transición prevista era la doble pila: añadir IPv6 sin quitar
IPv4 y apagar IPv4 más adelante. La primera mitad se ha cumplido, la segunda no.
El resultado es que el adoptante de IPv6 no ahorra nada: sigue necesitando su
parque de direcciones IPv4, sigue manteniendo la CGNAT y ahora además opera un
segundo protocolo, con su propio plan de direccionamiento, sus propias reglas de
cortafuegos y su propia formación. El coste es individual, inmediato y cierto; el
beneficio es colectivo, futuro e incierto. Es la definición de inercia excesiva
de \textcite{farrell1985}, y la conclusión a la que llega también el análisis
económico de la \textcite{oecd2014} y el estudio de adopción por actores de
\textcite{nikkhah2016}.

\section{La distribución que la media esconde}

\begin{figure}[ht]
  \centering
  \begin{tikzpicture}
    \foreach \xv/\lab in {0/0, 1.455/10, 2.910/20, 4.365/30, 5.820/40,
                          7.275/50} {
      \draw[gray!22] (\xv,-0.40) -- (\xv,3.40);
      \node[below, font=\tiny, gray!55!black] at (\xv,-0.40) {\lab\,\%};
    }
    \draw[red!65!black, thick, densely dashed] (6.294,-0.40) -- (6.294,3.40);
    \foreach \y/\x/\lab/\v in {%
      3.0/7.319/{Asia}/{50,30},
      2.4/7.222/{Américas}/{49,64},
      1.8/6.012/{Oceanía}/{41,32},
      1.2/5.110/{Europa}/{35,12},
      0.6/0.968/{África}/{6,65}} {
      \node[left, font=\scriptsize] at (-0.12,\y) {\lab};
      \fill[cIPv6] (0,\y-0.18) rectangle (\x,\y+0.18);
      \node[right, font=\scriptsize] at (7.45,\y) {\v\,\%};
    }
    \node[left, font=\scriptsize\bfseries] at (-0.12,0) {Mundo};
    \fill[cIPv6!45] (0,-0.18) rectangle (6.294,0.18);
    \node[right, font=\scriptsize\bfseries] at (7.45,0) {43,26\,\%};
  \end{tikzpicture}
  \caption{Usuarios con capacidad IPv6 por región, según la medición de APNIC
  Labs del 15 de agosto de 2026 \parencite{apnicv6}. La línea discontinua marca
  la media mundial. La distribución no es unimodal: Asia y las Américas han
  superado el 49 por ciento mientras África se queda en el 6,65, de modo que la
  media mundial no describe a ninguna región concreta. Elaboración propia.}
  \label{fig:regiones}
\end{figure}

La media del 43 por ciento oculta una distribución muy desigual, que además
invierte la intuición habitual: Asia y las Américas van por delante de Europa. La
explicación que da \textcite{huston2026b} es económica y no técnica. Las redes que
más han desplegado IPv6 son las móviles y las de incorporación reciente, como
Reliance Jio en la India, porque nacieron sin una base instalada de direcciones
IPv4 que amortizar. Quien ya tenía direcciones tenía un incentivo para
rentabilizarlas; quien no las tenía encontraba más barato empezar en IPv6. La
adopción no sigue al desarrollo económico, sigue a la ausencia de activos
heredados.

\section{¿Cuándo y por qué ocurrió?}

El Cuadro~\ref{tab:crono} responde al \enquote{cuándo} y el
Cuadro~\ref{tab:matriz} al \enquote{por qué}.

\begin{table}[ht]
  \centering
  \small
  \caption{Cronología verificada del proceso.}
  \label{tab:crono}
  \begin{tabularx}{\textwidth}{@{}l >{\raggedright\arraybackslash}X@{}}
    \toprule
    \textbf{Fecha} & \textbf{Hecho} \\
    \midrule
    09/1981 & RFC 791: especificación de IPv4, 32 bits de dirección
      \parencite{rfc791} \\
    11/1981 & RFC 801: plan de transición de NCP a TCP/IP, con fecha límite
      \parencite{rfc801} \\
    01/01/1983 & Día de la bandera: unos 400 servidores de ARPANET cambian de
      protocolo \\
    09/1993 & RFC 1519: CIDR elimina las clases de direcciones
      \parencite{rfc1519} \\
    05/1994 & \textbf{RFC 1631: NAT}, diecinueve meses antes que IPv6
      \parencite{rfc1631} \\
    12/1995 & RFC 1883: primera especificación de IPv6 \parencite{rfc1883} \\
    02/1996 & RFC 1918: bloques de direccionamiento privado
      \parencite{rfc1918} \\
    12/1998 & RFC 2460: segunda especificación de IPv6 \parencite{rfc2460} \\
    04/09/2008 & Google publica el primer dato de su serie: 0,14 por ciento
      \parencite{googlev6} \\
    03/02/2011 & IANA agota su reserva central de direcciones IPv4
      \parencite{nro2026} \\
    08/06/2011 & World IPv6 Day: prueba coordinada de 24 horas \\
    06/06/2012 & World IPv6 Launch: activación permanente \parencite{isoc2012} \\
    04/2012 & RFC 6598: bloque reservado para CGNAT \parencite{rfc6598} \\
    07/2017 & RFC 8200: IPv6 pasa a ser norma de Internet, STD 86
      \parencite{rfc8200} \\
    12/2021 & Máximo histórico del precio de una dirección IPv4, unos 55 dólares \\
    12/2025 & IPv6 cumple treinta años con menos de la mitad de los usuarios
      \parencite{register2025} \\
    28/03/2026 & \textbf{Primer día} en que la medición de Google supera el 50
      por ciento, con 50,10 \parencite{googlev6} \\
    12/08/2026 & Google mide 47,03 por ciento y APNIC Labs 43,16 por ciento
      \parencite{googlev6,apnicv6} \\
    \bottomrule
  \end{tabularx}
\end{table}

\begin{table}[ht]
  \centering
  \small
  \caption{Matriz de figuras de mérito de los dos protocolos.}
  \label{tab:matriz}
  \begin{tabularx}{\textwidth}{@{}>{\raggedright\arraybackslash}p{0.27\textwidth}
    >{\raggedright\arraybackslash}X >{\raggedright\arraybackslash}X
    >{\raggedright\arraybackslash}p{0.12\textwidth}@{}}
    \toprule
    \textbf{FOM} & \textbf{IPv4 (1981)} & \textbf{IPv6 (1995)}
      & \textbf{Ventaja} \\
    \midrule
    Espacio de direcciones
      & \num{4294967296} & $3{,}40\times10^{38}$ & IPv6, $\times 7{,}9\times10^{28}$ \\
    Direcciones por usuario de Internet
      & 0,70 & $5{,}6\times10^{28}$ & IPv6 \\
    Longitud de la cabecera
      & 20 a 60 octetos, variable & 40 octetos, fija & mixta \\
    Campos de la cabecera
      & 12 más opciones & 8 & IPv6 \\
    Suma de comprobación por salto
      & sí, se recalcula & no existe & IPv6 \\
    Fragmentación en tránsito
      & la hacen los routers & solo el origen & IPv6 \\
    Autoconfiguración sin servidor
      & no & sí (SLAAC) & IPv6 \\
    Entradas en la tabla global
      & \num{1116753} & \num{260375} & IPv6 \\
    Sobrecarga por paquete
      & 20 octetos mínimo & 40 octetos siempre & \textbf{IPv4} \\
    Tiempo de ida y vuelta
      & referencia & equivalente & ninguna \\
    Compatibilidad con la base instalada
      & total & \textbf{ninguna} & \textbf{IPv4} \\
    Coste de una dirección
      & unos 20 dólares & prácticamente cero & IPv6 \\
    Usuarios que lo pueden usar
      & 100 por ciento & 43 a 47 por ciento & \textbf{IPv4} \\
    \bottomrule
  \end{tabularx}
\end{table}

Cinco mecanismos explican el \enquote{por qué}, y ninguno de ellos es la calidad
del protocolo.

\textbf{01. La compatibilidad hacia atrás es una figura de mérito, y era la
única que importaba.} Un protocolo de red no vale por lo que hace, sino por con
cuántos permite hablar. En ese eje IPv6 salió al mercado con el valor cero y su
predecesor estaba en el cien por cien. Es el fallo que el propio IETF reconoció
en 2009 \parencite{daigle2009}.

\textbf{02. El beneficio es colectivo y el coste es individual.} Cada operador
que despliega IPv6 asume el gasto completo y no puede apagar IPv4, porque sus
clientes seguirían necesitando llegar a la mitad larga de los usuarios y de los
servidores que aún no tienen IPv6. Esperar es la respuesta racional de cada actor por separado y el
peor resultado para el conjunto \parencite{farrell1985,oecd2014,nikkhah2016}.

\textbf{03. El titular mejoró justo en el FOM que justificaba al invasor.} CIDR
en 1993, NAT en 1994, direccionamiento privado en 1996 y CGNAT en 2012
convirtieron un techo duro de \num{4294967296} direcciones en un techo elástico.
Este es el efecto vela en su forma más nítida, y con una precisión que los casos
clásicos no tienen: aquí el remiendo lleva número de norma y fecha de
publicación, y se adelantó diecinueve meses al sucesor.

\textbf{04. El invasor no ganaba en ningún FOM perceptible por el usuario.} Ni
más rápido, ni más seguro, ni más funcional \parencite{huston2024}. Una
tecnología cuya única ventaja es evitar un problema que el usuario no ha sufrido
todavía carece de argumento de venta.

\textbf{05. El único motor económico se está apagando.} El precio de una
dirección IPv4 cayó de unos 55 dólares a finales de 2021 a unos 20 en 2026
\parencite{huston2026a,ipv4center2026}. La escasez, que era la última fuerza
capaz de empujar la migración, se relaja en lugar de agravarse.

\section{Precisiones y limitaciones}

\begin{enumerate}
  \item \textbf{Las dos extrapolaciones de la Figura~\ref{fig:curvas} se
    contradicen y no puedo elegir entre ellas con estos datos.} El ajuste
    logístico da una asíntota del 47,6 por ciento con la serie de Google y del
    43,3 con la de APNIC Labs, es decir predice que la curva se detiene
    aproximadamente donde ya está. La regresión lineal de los últimos cinco años
    da 2,73 puntos porcentuales anuales y llega al cien por cien en 2045.
    Sobre esos mismos cinco años, la recta ajusta ligeramente mejor que la
    logística, con un error cuadrático medio de 1,81 puntos frente a 1,98. Además
    la logística presenta sesgo positivo en los dos últimos años, de $+0{,}92$ y
    $+1{,}42$ puntos, lo que indica que su asíntota está subestimada: ajustar una
    logística a datos que acaban de pasar el punto de inflexión sitúa siempre la
    asíntota cerca del valor actual. \textbf{La asíntota debe leerse como una
    advertencia, no como un pronóstico.} Lo robusto es lo otro: el ritmo anual
    cayó de 5,83 puntos en 2016 y 2017 a 2,41 en los dos últimos años, y a ese
    ritmo la transición necesitaría veinte años más, esto es, más de cincuenta
    desde la RFC 1883. Esa proyección coincide con la de \textcite{huston2024},
    que sitúa el final en 2045 con datos y método distintos.
  \item \textbf{Las tres mediciones públicas no coinciden y miden cosas
    distintas.} Google mide accesos de usuarios a sus servicios, APNIC Labs
    muestrea mediante anuncios publicitarios con ponderación por población y
    Cloudflare mide tráfico transportado. En agosto de 2026 dan 47, 43 y unos 40
    por ciento respectivamente. He usado las dos primeras porque publican la
    serie completa en bruto; las conclusiones no dependen de cuál se elija,
    porque la forma de las dos curvas es la misma.
  \item \textbf{El titular del 50 por ciento es un pico de fin de semana, no un
    nivel.} La serie de Google tiene un ciclo semanal marcado, con los sábados y
    domingos 2,92 puntos por encima de los días laborables, porque en casa se
    navega desde redes domésticas y móviles con más IPv6 que en las corporativas.
    De los \num{6552} días de la serie, solo 15 superan el 50 por ciento y todos
    son fines de semana; el 28 de marzo de 2026 fue sábado. La media móvil de 30
    días nunca ha llegado al 50: su máximo histórico es 48,49 por ciento, el 9 de
    agosto de 2026.
  \item \textbf{El tramo de IPv4 anterior a 1990 es esquemático.} No existe una
    serie continua del porcentaje de servidores migrados durante 1982 y 1983. La
    forma casi vertical de la curva refleja el hecho documentado de que la
    transición fue un cambio con fecha señalada \parencite{rfc801}, pero la
    pendiente exacta no está medida.
  \item \textbf{IPv6 no es un fracaso en términos absolutos.} Cuarenta y tres de
    cada cien usuarios de un universo de \num{6100} millones son unos
    \num{2600} millones de personas: es uno de los despliegues de infraestructura
    más grandes de la historia. Lo que ha fracasado es la \emph{sustitución}, no
    el despliegue. John Curran, presidente de ARIN, defiende justamente esa
    lectura: IPv6 \enquote{no iba de apagar IPv4, sino de asegurar que Internet
    pudiera seguir creciendo sin romperse} \parencite{register2025}.
  \item \textbf{No puedo contrastar el contrafactual.} La tesis de que la
    incompatibilidad hacia atrás fue la causa decisiva es coherente con la
    teoría, con los datos y con el reconocimiento explícito del IETF, pero no es
    comprobable: no existe una segunda Internet en la que se publicara un IPv6
    compatible y podamos medir si se adoptó más rápido.
  \item \textbf{El caso está abierto.} \textcite{huston2026a} señala que si la
    presión de escasez volviera a crecer lo suficiente, la base con IPv6 podría
    alcanzar el dominio de mercado y entonces el mercado de direcciones IPv4 se
    hundiría rápidamente. Los datos actuales apuntan en la dirección contraria,
    pero un cambio de régimen sigue siendo posible.
\end{enumerate}

\section*{Conclusiones}

IPv6 es, casi con seguridad, el mejor ejemplo disponible de avance técnicamente
superior que no se implementó a gran escala, y el enunciado admite una respuesta
precisa.

\textbf{Cuándo.} Nunca, y esa es la respuesta literal. En los \num{6552} días de
la serie de Google, IPv6 ha superado el 50 por ciento de los usuarios en quince,
todos ellos fines de semana, y la media móvil de 30 días nunca lo ha alcanzado.
Treinta años después de su especificación, y nueve después de convertirse en
norma de Internet, el sucesor no ha sustituido al predecesor en ningún momento
de su historia. El titular no ha cedido un solo punto: sigue en el cien por
cien.

\textbf{Por qué.} Porque IPv6 era superior en la figura de mérito del artefacto y
no en la figura de mérito del sistema. Ganaba por 29 órdenes de magnitud en
espacio de direcciones, un atributo que el usuario no percibe, y valía cero en
compatibilidad con la base instalada, el único atributo que gobierna la adopción
de un protocolo de red. Mientras tanto, el titular no se defendió compitiendo en
el terreno del invasor: mejoró en el suyo. CIDR, NAT y CGNAT convirtieron el
techo de direcciones de IPv4 de rígido en elástico, y lo hicieron antes de que
el sucesor existiera. Hoy, treinta años después, el precio de una dirección IPv4
baja en lugar de subir, que es la señal de mercado de que la escasez que
justificaba la migración ha dejado de apretar.

La lección que me llevo como CTO se puede formular en una frase:
\textbf{una tecnología que exige que todos se muevan a la vez no tiene una curva
de adopción, tiene un problema de coordinación}, y los problemas de coordinación
no se resuelven mejorando el producto. IPv6 dedicó treinta años a ser mejor
cuando lo que necesitaba era ser compatible. Y hay un corolario todavía más
incómodo para quien decide una hoja de ruta: \textbf{el paliativo que llega
primero fija el calendario del sucesor}. NAT se publicó diecinueve meses antes
que IPv6 y esos diecinueve meses de ventaja llevan tres décadas decidiendo el
resultado.

\printbibliography[title={Referencias bibliográficas}]

\end{document}
